\section{Hidden Meanings in Texts}

Since most people write in a stream of consciousness style, they are seldom aware that an esoteric text is deliberately constructed with several levels of meaning. The best example is Dante's description of the four levels of meaning in his texts.

\begin{enumerate}
\item \textbf{Literal}: The sense that does not go beyond the surface of the letter.
\item \textbf{Allegorical}: The truth hidden beneath a beautiful fiction.
\item \textbf{Moral}: The moral implications of a work of fiction.
\item \textbf{Anagogical}: The spiritual or mystical sense.
\end{enumerate}

The primary example is his \emph{Divine Comedy}. In brief:

\begin{enumerate}
\item Dante describes his journey through hell, purgatory, and heaven. There he encounters many people and sees many things.
\item He leaves many hints that reveal it is not a physical journey, but is a psychological and spiritual journey.
\item The sufferings in hell and purgatory reveal the moral meaning of the text.
\item The anagogical meaning is revealed in the Paradiso.
\end{enumerate}


A second type of esoteric meaning hides the true meaning for various reasons.

\begin{itemize}


\item The text includes information that should not be revealed to the public, who would just misunderstand it or misuse the information. Alchemical texts are in this class.

\item Esoteric teachings may be embodied in texts in order to protect them. Examples include fairy tales or myths. The public repeats them, without understanding them. However, only a few, possibly generations later, will be able to discern the true meanings.

\item The true teachings are a danger to the public order, so they are hidden in innocuous texts. Only a few will be able to unravel the true meanings.

\item The true teachings are dangerous, because the authorities may punish the author.
\end{itemize}


Gornahoor has secret meanings which some understand better than others. Some things are more deeply hidden. Perhaps a professor, 25 years from now, will try to unravel them. In the meantime, here are some simpler examples just for edification.

\paragraph{Double Chocolate Chip Muffins}


At the literal level, Double Chocolate Chip Muffins\footnote{\url{https://gornahoor.net/?p=14871}} is a recipe for tasty muffins. Most readers will stop there. But the recipe cannot be understood without actually making and then eating the muffins.

The same principle applies to other posts about self-knowledge, self-mastery, meditation, purification, and the like. You cannot just read about it. Rather, you must actually meditate, purify, etc. It amazes me how few people really grasp that. In almost all cases, it is necessary to find someone who can show you the way.

\paragraph{Hunting Iguanas}


At the literal level, Hunting Iguanas\footnote{\url{https://gornahoor.net/?p=14935}} is about ridding the environment of pests. But on the allegorical lever, it is about mastering one's thoughts. The lizards represent random thoughts and desires that upset the environmental balance of your consciousness. One must then be diligent to notices them and remove them. Otherwise, they will take over your psychic energy.

At the moral level, the lizards are amoral. Their attitude toward property is ``see and take'', and they have no regard for property rights. Perhaps, there is a mystical lizard experience; you tell me.

\paragraph{The Mathematician, the Carpenter, and the Physicist}


The Carpenter represents the Literal level\footnote{\url{https://gornahoor.net/?p=14883}}, since he is interested in solving a physical problem. They physicist is at the allegorical level. Physicists re-interpret the math equations as forces, etc., even when there is no physical reality that corresponds to the mathematical equations. Mathematically, there is no direction of time nor even of motion.

The mathematician is at the moral level and the bartender at the mystical level.


\flrightit{Posted on 2021-10-14 by Cologero}

\begin{center}* * *\end{center}

\begin{footnotesize}
\begin{sffamily}
\texttt{Greg on 2021-10-15 at 15:04 said:}

The Iguana represent the thoughts and desires of the lower intellectual and lower emotional centers. The lower intellect analytical intellect cannot reach God because He is separated from it by a cloud of unknowing. That is God is incomprehensible and transcendent and cannot be reached by the lower centers. Just as the Iguana cannot reach the meat in the trap so the lower thoughts, emotions and desires cannot reach God. When the thoughts drop away then the higher intellectual centers in union with the higher eros can know God in ;a unitive and intuitive manner. This is what is meant by infused contemplation. In the same way when the lower thoughts and desires fall away which Tomberg called concentration without effort, the Hermeticist is able to know the inner meaning of Symbols.

\hfill

\texttt{Niko The Exile on 2021-11-03 at 10:53 said:}

So this Cologero persona claims to actually know Language of the birds (or Phonetic cabala)


\hfill

\end{sffamily}
\end{footnotesize}

